\documentclass[12pt]{article}
\usepackage[utf8]{inputenc}
\usepackage[T1]{fontenc}
\usepackage{geometry}
\usepackage{setspace}
\usepackage{booktabs}
\usepackage{longtable}
\usepackage{array}
\usepackage{enumitem}
\usepackage{hyperref}
\usepackage{graphicx}
\usepackage{listings}
\usepackage{color}
\geometry{margin=1in}
\setlength{\parskip}{0.55em}
\setlength{\parindent}{0pt}
\onehalfspacing

\definecolor{listinggray}{gray}{0.95}
\lstset{
  basicstyle=\ttfamily\small,
  backgroundcolor=\color{listinggray},
  frame=single,
  breaklines=true,
  columns=fullflexible
}

\begin{document}

\begin{titlepage}
    \newgeometry{margin=0pt}
    \noindent\includegraphics[width=\paperwidth,height=\paperheight]{pagedegardeNLP.png}
    \restoregeometry
\end{titlepage}

\begin{titlepage}
    \centering
    {\Huge Final Report\\TuniSpeak RAG Assistant\\[0.35cm]}
    {\Large CRISP-DM Project Overview}\vfill
    {\large TuniSpeak Team\\December 2025}\vfill
    {\footnotesize Internal document -- restricted circulation}\vspace*{0.6cm}
\end{titlepage}

\thispagestyle{empty}
\tableofcontents
\newpage

\section*{Executive Summary}
TuniSpeak is a trilingual (French, Modern Standard Arabic, Tunisian Darija) question-answering assistant for Tunisian universities. The system combines a FastAPI backend, a Streamlit comparison dashboard, a hybrid Retrieval-Augmented Generation (RAG) pipeline, and an optional voice interface. This report follows the CRISP-DM methodology and reflects the state of repository \texttt{TuniSpeak} (branch \texttt{main}) as of 9 December 2025.

\begin{itemize}
    \item FAQ corpora in \texttt{data/faq/} produce 3\,343 curated knowledge chunks (\texttt{data/processed/chunks.jsonl}).
    \item Hybrid retrieval fuses BM25 with multilingual sentence-transformer embeddings and an optional CrossEncoder reranker.
    \item Answers originate from \texttt{deepset/xlm-roberta-base-squad2}; TinyLlama LoRA adapters polish responses when enabled.
    \item The Streamlit UI exposes Bronze/Silver/Gold/Platinum modes to compare retrieval depth, metadata, feedback, and voice capabilities.
    \item Latest automated evaluation (\texttt{data/processed/eval/evaluation\_summary.json}) reports EM 0.479, F1 0.674, Recall@5 0.969, nDCG@5 0.767, abstention 0.021, and ECE 0.127.
\end{itemize}

\section{Business Understanding}
\subsection{Stakeholders and Objectives}
Key stakeholders include student services teams, university leadership, and enrolled students. The project aims to deliver 24/7 multilingual guidance, reduce repetitive ticket volume, and log feedback to drive continuous improvement.

\subsection{Pain Points}
Help desks experience seasonal backlogs, particularly during enrolment and scholarship periods. Darija queries often lack written resources, forcing manual support. Students expect immediate, reliable answers with clear provenance.

\subsection{Success Criteria}
\begin{itemize}
    \item Production endpoint availability via \texttt{poetry run tunispeak-api} under campus infrastructure constraints.
    \item Recall@5 \geq 0.95 while keeping abstention below 3\% for risk-aware responses.
    \item Comprehensive telemetry in \texttt{data/processed/eval/history.jsonl} and feedback capture through \texttt{/api/v1/feedback}.
    \item Ability to audit responses with citations and confidence scores across all UI modes.
\end{itemize}

\section{Data Understanding}
\subsection{Source Inventory}
\begin{itemize}
    \item \texttt{data/faq/dataset\_finetune.jsonl}: 1\,117 dialogues (434 French, 372 Arabic, 311 Darija/Arabizi).
    \item \texttt{data/faq/mini\_faq.jsonl}: 48 curated QA pairs for rapid evaluation cycles.
    \item \texttt{data/raw/}: official documents (PDF, DOCX, TXT) ingested during initial runs.
    \item User feedback: \texttt{data/feedback/feedback.jsonl} storing helpful/unhelpful votes and suggested revisions.
\end{itemize}

\subsection{Ingestion Overview}
The Typer CLI \texttt{poetry run tunispeak-ingest run <folder>} triggers \texttt{DocumentIngestionService}. Supported formats include \texttt{.jsonl}, \texttt{.json}, \texttt{.pdf}, \texttt{.docx}, \texttt{.txt}, and \texttt{.md}. Each document is segmented, language-detected, and written to \texttt{data/processed/chunks.jsonl}. Canonical language mapping collapses Tunisian Darija variants to the ISO code \texttt{aeb}, ensuring consistent downstream processing.

\begin{lstlisting}[language=Python,caption={Chunk staging inside DocumentIngestionService}]
def _stage_chunk(self, path: Path, index: int, chunk_text: str) -> None:
    chunk_text = chunk_text.strip()
    if not chunk_text:
        return
    detection = self.detector.detect(chunk_text)
    record = {
        "document_id": path.stem,
        "chunk_id": f"{path.stem}-{index:04d}",
        "text": chunk_text,
        "language": detection.language,
        "normalized_text": detection.normalized,
        "source_path": str(path.resolve()),
    }
    self.buffer.append(record)
\end{lstlisting}

\subsection{Language Detection}
\texttt{LanguageDetector} in \texttt{app/services/language.py} combines script heuristics, Arabizi markers, and \texttt{langdetect}. New tokens such as \textit{chnoa}, \textit{bach}, and \textit{lazimni} resolve previous 422 errors for romanised Darija queries.

\subsection{Dataset Snapshot}
\begin{longtable}{>{\raggedright\arraybackslash}p{5cm} >{\raggedright\arraybackslash}p{9cm}}
\toprule
Statistic & Value \\
\midrule
Chunks in \texttt{data/processed/chunks.jsonl} & 3\,343 \\
Language distribution & \texttt{aeb}: 2\,914; \texttt{fr}: 426; \texttt{ar}: 2; other: 1 \\
Principal documents & \texttt{dataset\_finetune} (1\,117); \texttt{dataset\_final\_universite\_tunisie} (1\,082); \texttt{mini\_faq} (48) \\
\bottomrule
\end{longtable}

\section{Data Preparation}
\subsection{Cleaning and Normalisation}
Ingestion trims whitespace, removes empty segments, normalises Darija tokens, and standardises language codes. For romanised Darija, \texttt{normalize\_darija} unifies common spellings (e.g., \texttt{ch} \textrightarrow{} \texttt{sh}, \texttt{bech} \textrightarrow{} \texttt{bash}). Each chunk stores both raw and normalised text, enabling search over consistent representations.

\subsection{Quality Controls}
\begin{itemize}
    \item Empty or non-text documents emit warnings for operator review.
    \item Flushing rehydrates chunks to keep ingestion idempotent across runs.
    \item Telemetry and feedback rely on JSON Lines for easy streaming into dashboards.
\end{itemize}

\subsection{Feature Preparation}
\texttt{scripts/build\_hybrid\_index.py} builds BM25 statistics and dense embeddings. The dense encoder defaults to \texttt{sentence-transformers/paraphrase-multilingual-MiniLM-L12-v2}; embeddings are persisted alongside a pickled BM25 model for reproducibility.

\section{Modeling}
\subsection{RAG Pipeline}
\texttt{QAPipeline} orchestrates four stages:
\begin{enumerate}[label=\textbf{\alph*)}]
    \item Language detection to normalise the query and bias retrieval toward the requester's language.
    \item Hybrid retrieval combining BM25 and dense cosine scores with optional CrossEncoder reranking.
    \item Extractive QA powered by \texttt{deepset/xlm-roberta-base-squad2} with language-specific confidence thresholds and abstention overrides.
    \item Optional answer synthesis via TinyLlama LoRA adapters to improve fluency while preserving citations.
\end{enumerate}

\subsection{Complexity Modes}
The Streamlit UI queries the same API endpoint with differing parameters:
\begin{itemize}
    \item Bronze: \texttt{top\_k=3}, no metrics or sources.
    \item Silver: \texttt{top\_k=5}, displays detected language, confidence, and up to three citations.
    \item Gold: \texttt{top\_k=7}, adds the normalised query and in-app feedback form.
    \item Platinum: \texttt{top\_k=10}, removes citation limits and unlocks speech (WebRTC + Faster-Whisper + gTTS).
\end{itemize}

\subsection{Speech Stack}
\texttt{SpeechService} leverages Faster-Whisper (model size ``small'') for transcription and gTTS for synthesis. The \texttt{/api/v1/qa/voice} route returns both the textual answer and base64-encoded MP3 audio for playback in the Platinum UI.

\section{Evaluation}
\subsection{Aggregate Metrics}
\begin{table}[h]
    \centering
    \begin{tabular}{lc}
        \toprule
        Indicator & Value \\
        \midrule
        Exact Match (EM) & 0.479 \\
        F1 & 0.674 \\
        Recall@5 & 0.969 \\
        nDCG@5 & 0.767 \\
        Abstention rate & 0.021 \\
        Expected Calibration Error (ECE) & 0.127 \\
        Brier score & 0.993 \\
        \bottomrule
    \end{tabular}
    \caption{Extracted from \texttt{data/processed/eval/evaluation\_summary.json}}
\end{table}

High recall confirms the hybrid retriever surfaces relevant evidence. Lower EM indicates paraphrasing differences and limited Darija coverage. ECE above 0.1 highlights calibration improvements as a follow-up action.

\subsection{Representative Error}
\begin{itemize}
    \item Query: ``شنية الاجراءات باش نبدل الشعبة متاعي؟''
    \item Prediction: generic guidance lacking document-specific steps.
    \item Mitigation: extend \texttt{dataset\_final\_universite\_tunisie.json} with programme-change procedures.
\end{itemize}

\section{Deployment}
\subsection{Runtime Components}
\begin{itemize}
    \item FastAPI + Uvicorn launched via \texttt{poetry run tunispeak-api}.
    \item Streamlit UI via \texttt{poetry run streamlit run streamlit\_app.py}.
    \item Telemetry logging in \texttt{data/processed/eval/history.jsonl} with \texttt{TelemetryLogger}.
    \item Feedback persisted by \texttt{/api/v1/feedback} into \texttt{data/feedback/feedback.jsonl}.
\end{itemize}

\subsection{Operational Scripts}
\begin{longtable}{>{\raggedright\arraybackslash}p{4cm} >{\raggedright\arraybackslash}p{10cm}}
\toprule
Command & Purpose \\
\midrule
\texttt{poetry run tunispeak-ingest run data/faq} & Refresh chunks from FAQ corpora. \\
\texttt{poetry run tunispeak-index} & Build BM25 and dense indices. \\
\texttt{poetry run tunispeak-eval --dataset data/faq/mini\_faq.jsonl} & Produce evaluation summary. \\
\texttt{poetry run tunispeak-api} & Start the FastAPI service. \\
\texttt{poetry run streamlit run streamlit\_app.py} & Launch the comparison interface. \\
\bottomrule
\end{longtable}

\subsection{Limitations and Next Steps}
\begin{itemize}
    \item Darija coverage remains skewed: only two Arabic-script chunks and 426 French chunks versus 2\,914 Darija entries; prioritise additional Arabic sources.
    \item Calibration requires temperature scaling or isotonic regression to reduce ECE below 0.1.
    \item Monitoring lacks Prometheus or OpenTelemetry exporters for production observability.
    \item CI/CD pipelines are absent; integrate tests and lint checks into GitHub Actions.
\end{itemize}

\section{Conclusion}
TuniSpeak delivers a targeted RAG assistant: hybrid retrieval, multilingual normalisation, evidence-backed answers, and a progressive UI with speech support. The system already offers strong recall and logging foundations. Future work should expand Darija resources, improve calibration, and add observability tooling for production deployment.

\appendix
\section{Key Configuration (excerpt from \texttt{app/core/config.py})}
\begin{lstlisting}[language=Python]
class Settings(BaseSettings):
    dense_model = "sentence-transformers/paraphrase-multilingual-MiniLM-L12-v2"
    cross_encoder_model = "cross-encoder/ms-marco-MiniLM-L-6-v2"
    qa_model = "deepset/xlm-roberta-base-squad2"
    qa_confidence_threshold = 0.15
    qa_language_thresholds = {"fr": 0.05, "ar": 0.02, "aeb": 0.005}
    speech_model_size = "small"
    speech_enable_tts = True
    telemetry_log_path = "./data/processed/eval/history.jsonl"
    feedback_log_path = "./data/feedback/feedback.jsonl"
\end{lstlisting}

\section{Useful Commands}
\begin{itemize}
    \item \texttt{poetry run python scripts/ingest\_documents.py run data/raw}
    \item \texttt{poetry run python scripts/build\_hybrid\_index.py}
    \item \texttt{poetry run pytest}
\end{itemize}

\section{Glossary}
\begin{description}[style=nextline]
    \item[CRISP-DM] Cross-Industry Standard Process for Data Mining; framework with phases Business Understanding, Data Understanding, Data Preparation, Modeling, Evaluation, Deployment.
    \item[EM] Exact Match; proportion of answers identical to the reference.
    \item[ECE] Expected Calibration Error; discrepancy between predicted confidence and empirical accuracy.
    \item[RAG] Retrieval-Augmented Generation; combines document retrieval with generation.
    \item[LoRA] Low-Rank Adaptation; parameter-efficient fine-tuning method for large models.
\end{description}

\end{document}
